% !PRINTED 45
\weberpage{45}
%<<CONT

3. \gesperrt{Die Werthe eines Charakters $\chi(A)$ sind Einheitswurzeln, deren Grad ein Theiler des Grades von $A$ ist.}

Danach können wir leicht alle Charaktere bestimmen. Setzen wir nämlich
\weq{9}{
\chi(A_1) = \omega_1,\ \chi(A_2) = \omega_2 \ldots \chi(A_\nu) = \omega_\nu,
}
so ist
\weq{10}{\omega_1^{a_1} = 1,\ \omega_2^{a_2} = 1, \ldots \omega_\nu^{a_\nu} = 1,}
und es ist nach (2) und (4)
\weq{11}{\chi(A) = \omega_1^{\alpha_1}\omega_2^{\alpha_2}\ldots\omega_\nu^{\alpha_\nu}.}

Umgekehrt ist, wenn $\omega_1, \omega_2 \ldots \omega_\nu$ irgend eine Lösung der Gleichungen (10) bedeutet, durch (11) eine Function von $A$ bestimmt, die nach 1. der Bedingung (4) genügt und also ein Charakter der Gruppe ist.

Jede der Gleichungen (10) hat $a_1, a_2 \ldots a_\nu$ Wurzeln, und wenn wir jede mit jeder combiniren, so ergeben sich
\[
a_1 a_2 \ldots a_\nu = n
\]
Combinationen. Alle diese Combinationen führen zu verschiedenen Charakteren $\chi(A)$. Denn bezeichnen wir mit $\omega_1', \omega_2' \ldots \omega_\nu'$ eine zweite Combination von Wurzeln der Gleichung (10), und ist für alle Elemente $A$
\[
\omega_1^{\alpha_1}\omega_2^{\alpha_2}\ldots\omega_\nu^{\alpha_\nu} = \omega_1'^{\alpha_1}\omega_2'^{\alpha_2}\ldots\omega_\nu'^{\alpha_\nu},
\]
so folgt, wenn man $\alpha_1 = 1, \alpha_2 = 0 \ldots \alpha_\nu = 0$ setzt, $\omega_1 = \omega_1'$ und ebenso $\omega_2 = \omega_2' \ldots \omega_\nu = \omega_\nu'$. Daraus folgt der Satz:

4. \gesperrt{Es giebt $n$ und nicht mehr verschiedene Charaktere einer Abel'schen Gruppe $n^{\text{ten}}$ Grades, die alle durch die Formel (11) dargestellt sind.}

Wenn wir unter $\varepsilon_1, \varepsilon_2 \ldots \varepsilon_\nu$ ein System primitiver Wurzeln der Gleichungen (10) verstehen, so können wir
\weq{12}{\omega_1 = \varepsilon_1^{\beta_1},\ \omega_2 = \varepsilon_2^{\beta_2}, \ldots \omega_\nu = \varepsilon_\nu^{\beta_\nu}}
setzen, und können die $\beta_1, \beta_2 \ldots \beta_\nu$ ebenso wie die $\alpha_1, \alpha_2 \ldots \alpha_\nu$ je einem vollen Restsysteme nach den Moduln $a_1, a_2 \ldots a_\nu$ entnehmen. Dann bekommen wir die sämmtlichen $n$ Gruppencharaktere bei feststehenden $\varepsilon_1, \varepsilon_2 \ldots \varepsilon_\nu$ in der Form
\weq{13}{\chi(A) = \varepsilon_1^{\alpha_1\beta_1}\varepsilon_2^{\alpha_2\beta_2}\ldots\varepsilon_\nu^{\alpha_\nu\beta_\nu}.}
